% ==================== 第 8 章 总结与展望 ====================
\chapter{总结与展望}
\label{ch:conclusion}

\section{工作总结}

本文围绕“在 Host 不可信前提下保护密钥敏感操作”这一核心问题，设计并实现了一套 Host 与 Enclave 分离的轻量级密钥管理系统，主要工作包括：

\begin{enumerate}
  \item \textbf{架构与密钥层次。} 提出并实现了 Host + Enclave 分离的三层架构，以信封加密构建“主密钥—数据密钥—业务数据”的层次结构，使主密钥仅驻留于模拟 Enclave，数据库仅保存包装后的 DEK 与业务密文。
  \item \textbf{协议实现。} 实现了 Host 与 Enclave 之间的远程认证握手与基于 AES-256-GCM、带序列号防重放的安全信道，并在 Enclave 内完成 DEK 的包装/解包与数据加解密，使明文 DEK 不出 Enclave。
  \item \textbf{系统功能。} 实现了 JWT 鉴权、基于角色的访问控制、版本化密钥轮换、密钥全生命周期管理与审计日志，并提供本地与可信两种可对照的运行模式。
  \item \textbf{威胁模型与验证。} 建立了系统的威胁模型与安全需求（SR1--SR8），逐项给出已实现的缓解机制，并诚实界定了软件模拟 TEE 的能力边界；以 8 项自动化集成测试验证了系统在两种模式下的功能正确性，并提供一键启动以保证可复现。
\end{enumerate}

综上，在“Host 不可信、Enclave 相对可信”的受控威胁模型下，本文方案能够将主密钥与数据密钥的敏感操作隔离至 Enclave，达成了既定研究目标。

\section{局限性}

本着诚实表述的原则，本文系统存在以下局限，这些局限与第~\ref{ch:threat}~章的边界声明一致：

\begin{itemize}
  \item \textbf{软件模拟而非硬件隔离。} Enclave 为普通操作系统进程，主密钥驻留于普通进程内存，不具备硬件内存加密与隔离，无法抵御物理攻击、内存转储与 Enclave 自身被攻破等威胁。
  \item \textbf{远程认证为结构性模拟。} 身份密钥为对称密钥且握手阶段以明文传输，缺乏硬件根信任，无法抵御对握手过程的主动中间人攻击，其有效性依赖本地可信握手信道这一假设。
  \item \textbf{授权粒度与多租户。} 访问控制为角色级，任意已鉴权用户可使用任意 active 密钥，缺乏逐密钥 ACL 与租户隔离。
  \item \textbf{工程完整度。} 仍使用 SQLite、单 Enclave 实例、无高可用集群；容器化仅覆盖 Host 侧；未实现传输层 TLS（依赖部署层）、令牌吊销与自动轮换策略。
  \item \textbf{未做形式化验证。} 远程认证与信道协议未经过形式化建模与证明，性能也尚未给出实测数据。
\end{itemize}

\section{未来工作}

针对上述局限，后续可在以下方向展开：

\begin{enumerate}
  \item \textbf{对接真实硬件 TEE。} 将模拟 Enclave 替换为 Intel SGX 或 AWS Nitro Enclaves，采用硬件根信任的非对称远程认证，从根本上提升隔离强度。
  \item \textbf{后量子密钥封装。} 引入 ML-KEM 等后量子算法\cite{fips203}，对 DEK 的封装采用经典与后量子混合方案，提升前瞻安全性。
  \item \textbf{工程化增强。} 迁移至 PostgreSQL 并引入数据库迁移；实现密钥审批流、自动轮换策略、令牌吊销与多租户隔离；将 Enclave 独立容器化部署。
  \item \textbf{形式化与性能。} 对远程认证与安全信道协议进行形式化建模与分析；补充本地与可信模式的性能对比实测，量化 Enclave IPC 开销。
  \item \textbf{管理与演示界面。} 提供管理后台与可视化演示，提升系统可用性与展示效果。
\end{enumerate}

总体而言，本文完成了一个定位清晰、边界诚实、可复现的应用密码学工程原型，为后续对接真实 TEE 与后量子密钥封装提供了可扩展的基础。
